\documentclass[11pt,a4paper]{article}
\usepackage{curso}
\begin{document}
\portada{Ejercicios --- Lección}{9.3}{Modelos termodinámicos de eliminación de ruido}
\section{Ejercicios}
\begin{enumerate}[label=\textbf{E\arabic*.}]
\item (B) ¿Por qué un muestreador local mezcla lento entre modos separados?
\item (B) Escriba la energía de un condicional en una DTM gaussiana con núcleo $q(x_t\mid x_{t-1})=\N(cx_{t-1},1-c^2)$.
\item (B) Calcule $E$ para $T=4$, $K_{\text{mix}}=250$, $L=70$, $E_{\text{cell}}=2$ fJ.
\item (B) Enumere tres advertencias sobre la estimación de $10^4\times$.
\item (I) Con una barrera $\Delta E=20$ (en $k_BT$), estime el factor de Arrhenius de los cruces y compárelo con 8000 pasos.
\item (I) ¿Por qué con muy pocos pasos de ruido los condicionales vuelven a ser multimodales?
\item (A) Reproduzca el experimento de la lección con mallas de tiempo geométricas y estudie la fracción en el modo derecho frente al número de niveles.
\end{enumerate}
\section{Soluciones}
\begin{solucion}
\textbf{E1.} La probabilidad de subir $\Delta E$ es $\sim e^{-\Delta E}$: el tiempo entre cruces crece exponencialmente.

\textbf{E2.} $E(x_{t-1})+\frac{(x_t-cx_{t-1})^2}{2(1-c^2)}$, con $E=-\log p_{t-1}$.

\textbf{E3.} $4\cdot250\cdot4900\cdot2\ \text{fJ}\approx9{,}8$ nJ.

\textbf{E4.} Es una estimación (no un chip completo); energía de GPU subestimada por FLOPs/especificación; banco pequeño.

\textbf{E5.} $e^{-20}\approx2\times10^{-9}$: ni siquiera un cruce esperado en 8000 pasos.

\textbf{E6.} Si $q$ añade mucho ruido en un paso, $x_t$ es compatible con varios modos de $p_{t-1}$: el posterior hereda la multimodalidad.

\textbf{E7.} Compare con la malla uniforme ($75\,\%$ con 32 niveles); interprete el resultado en términos de dónde se separan los modos.
\end{solucion}
\end{document}
